% 两变量能量的 s-t 割表示。移位后一元容量：s->1 为 2（=phi~_1(1)），1->t 为 0；s->2 为 0，2->t 为 3（=phi~_2(0)）。
% 成对异值惩罚 w=4：1->2、2->1 各一条容量4的弧。源侧标签0、汇侧标签1，割容量 = J(x)+1。
% (a) 源侧 {s}：割断 s->1，容量2，对应 x=(1,1)，J=1（最小割）。(b) 源侧 {s,1}：割断 1->2，容量4，对应 x=(0,1)，J=3。
\begin{tikzpicture}[line width=.8pt,font=\figurefont\fontsize{8}{10}\selectfont,>=Stealth,
  vtx/.style={circle,draw=FigureStroke,fill=FigureBlueFill,minimum size=7mm,inner sep=0pt,line width=.8pt},
  capl/.style={font=\figurefont\fontsize{8}{10}\selectfont,inner sep=1.5pt},
  zero/.style={->,FigureGrid,dotted,line width=.8pt}]
  % (a)
  \node[vtx,fill=white] (as) at (0,0) {$s$};
  \node[vtx] (a1) at (1.9,1.0) {$1$};
  \node[vtx] (a2) at (1.9,-1.0) {$2$};
  \node[vtx,fill=white] (at) at (3.8,0) {$t$};
  \draw[zero] (a1)--node[capl,above right,FigureMuted]{0}(at);
  \draw[zero] (as)--node[capl,below left,FigureMuted]{0}(a2);
  \draw[->,FigureRed,line width=1.3pt] (as)--node[capl,pos=.62,above left,FigureInk]{2}(a1);
  \draw[->,FigureStroke] (a2)--node[capl,below right]{3}(at);
  \draw[->,FigureStroke] (a1) to[bend left=22] node[capl,right]{4}(a2);
  \draw[->,FigureStroke] (a2) to[bend left=22] node[capl,left]{4}(a1);
  \draw[FigureRed,dashed,line width=.8pt] (0.8,1.7)--(0.8,-1.7);
  \node[FigureInk,align=center,anchor=north] at (1.9,-1.8) {(a) 源侧$\{s\}$，$\bm x=(1,1)$\\割容量$2=J(\bm x)+1$};
  % (b)
  \begin{scope}[xshift=58mm]
  \node[vtx,fill=white] (bs) at (0,0) {$s$};
  \node[vtx] (b1) at (1.9,1.0) {$1$};
  \node[vtx] (b2) at (1.9,-1.0) {$2$};
  \node[vtx,fill=white] (bt) at (3.8,0) {$t$};
  \draw[zero] (b1)--node[capl,above right,FigureMuted]{0}(bt);
  \draw[zero] (bs)--node[capl,below left,FigureMuted]{0}(b2);
  \draw[->,FigureStroke] (bs)--node[capl,above left]{2}(b1);
  \draw[->,FigureStroke] (b2)--node[capl,below right]{3}(bt);
  \draw[->,FigureRed,line width=1.3pt] (b1) to[bend left=22] node[capl,right,FigureInk]{4}(b2);
  \draw[->,FigureStroke] (b2) to[bend left=22] node[capl,left]{4}(b1);
  \draw[FigureRed,dashed,line width=.8pt] plot[smooth] coordinates {(0.2,-1.6) (1.0,-0.3) (2.7,-0.1) (3.1,1.6)};
  \node[FigureInk,align=center,anchor=north] at (1.9,-1.8) {(b) 源侧$\{s,1\}$，$\bm x=(0,1)$\\割容量$4=J(\bm x)+1$};
  \end{scope}
\end{tikzpicture}
